%% ch05-operations.tex — JA4proxy Operator's User Guide
%% Chapter 5: Day-to-Day Operations

\chapter{Day-to-Day Operations}
\label{ch:operations}
\index{operations}

\newthought{Running \ja\ day-to-day is mostly unremarkable.} The proxy handles
connections, Redis holds state, and Grafana shows you what is happening. This
chapter covers the routine operational tasks: starting and stopping the stack,
updating lists, managing bans, maintaining the GeoIP database, and keeping
the proxy healthy.

\section{Starting and Stopping}
\index{make start}
\index{make stop}

\subsection{Normal Start}

\begin{bashcode}
make start
\end{bashcode}

Starts all containers. Existing Redis data (bans, rate-limit windows) is preserved.
The proxy begins processing connections as soon as it becomes healthy, which
takes approximately 5--10 seconds after the containers start.

\subsection{Stop Without Data Loss}

\begin{bashcode}
make stop
\end{bashcode}

Stops all containers. Redis data persists on the host volume. Logs are preserved.
Use this for planned maintenance or before deploying a configuration change
that requires a restart.

\subsection{Stop and Clean Everything}

\begin{bashcode}
make stop-clean
\end{bashcode}

Stops all containers and removes volumes. Use this when you want a completely
fresh start — for example, when moving from testing to a fresh production deployment,
or when debugging requires a known-clean Redis state.

\begin{dangerbox}[Data loss]
\code{make stop-clean} permanently deletes all Redis data: bans, rate-limit windows,
beaconing history, enrichment cache, and return visitor records. Any active bans
will be lifted immediately.
\end{dangerbox}

\subsection{Full Rebuild}

After pulling a new version of the repository:

\begin{bashcode}
git pull
make rebuild
\end{bashcode}

This stops the stack, rebuilds all Docker images from scratch, and starts the
new version.

\section{Updating Fingerprint Lists}
\index{fingerprint lists!updating}
\index{hot reload}

The JA4 blacklist, whitelist, patterns, and labels all live in
\code{config/proxy.yml} under \configkey{security}. To add or remove entries:

\begin{enumerate}
  \item Edit \code{config/proxy.yml}
  \item Reload without restart:
\begin{bashcode}
docker compose kill -s HUP ja4proxy
\end{bashcode}
  \item Verify the reload in logs:
\begin{bashcode}
docker compose logs ja4proxy | tail -5
# Should show: INFO | config | event=reload | status=ok
\end{bashcode}
\end{enumerate}

\begin{notebox}[Reloading in Production / Native Mode]
If JA4proxy is running in production or native systemd mode, trigger a configuration reload by sending \code{SIGHUP} directly to the Go daemon:
\begin{bashcode}
sudo kill -HUP $(pgrep ja4pd)
\end{bashcode}
To view logs and verify the reload:
\begin{bashcode}
sudo journalctl -u ja4proxy.service -n 5
# Should show: INFO | config | event=reload | status=ok
\end{bashcode}
\end{notebox}

\subsection{Adding a Fingerprint to the Blacklist}

When you identify a malicious fingerprint (from logs, Grafana, or threat intelligence):

\begin{yamlcode}
security:
  blacklist:
    - t13d0912_a1b2c3d4e5f6_987654321abc  # Sliver C2 — added 2026-04-04
    - t12d1302_c0ffee123456_aabbccddeeff  # Cobalt Strike — added 2026-04-04
    # Add your new entry here:
    - t13d1015_newfingerprint_0123456789  # Description — date
\end{yamlcode}

Comment each entry with a description and the date it was added. This is your
operational audit trail.

\subsection{Adding a Fingerprint to the Whitelist}

For a legitimate tool that is scoring too high:

\begin{yamlcode}
security:
  whitelist:
    - t13d1113h2_8daaf6152771_b0da82dd1658  # Chrome 120 (Windows)
    # Add your tool here:
    - t13d0914h2_abcdef123456_fedcba654321  # Internal monitoring agent
\end{yamlcode}

\begin{warningbox}[Whitelist changes are immediate]
Whitelist additions take effect on reload and apply to all subsequent connections.
They do not apply retroactively to active bans — if a whitelisted IP is currently
banned, the ban persists until it expires (5 minutes).
\end{warningbox}

\section{Managing IP Bans}
\index{bans!managing}
\index{Redis!ban keys}

\subsection{Viewing Current Bans}

\begin{bashcode}
# List all active bans with TTLs
docker compose exec redis redis-cli --scan --pattern "ban:*" | \
  xargs -I{} sh -c 'echo -n "{}: "; docker compose exec redis redis-cli TTL "{}"'
\end{bashcode}

Or for a quick count:

\begin{bashcode}
docker compose exec redis redis-cli --scan --pattern "ban:*" | wc -l
\end{bashcode}

\subsection{Clearing a Specific Ban}

If a legitimate IP has been incorrectly banned:

\begin{bashcode}
# Replace with the actual IP address
docker compose exec redis redis-cli DEL "ban:192.0.2.1"
\end{bashcode}

The next connection from that IP will be processed normally. There is no need
to restart anything.

\begin{notebox}[Redis Commands in Native Mode]
If you are running in native mode, Redis runs directly on the host. You can use \code{redis-cli} commands directly on the host terminal instead of wrapping them in \code{docker compose exec}:
\begin{bashcode}
# Count active bans directly
redis-cli --scan --pattern "ban:*" | wc -l

# Clear a specific ban directly
redis-cli DEL "ban:192.0.2.1"
\end{bashcode}
\end{notebox}

\subsection{Ban Expiry}

All bans have a 5-minute TTL.\sidenote{The self-healing TTL is intentional and cannot
be disabled. See Section~1.5 for the rationale.} You do not need to manually clear
bans — they expire automatically. Manual clearing is only needed when you need to
restore access immediately (e.g., during an incident where a legitimate user has been
blocked).

\subsection{Clearing All Bans}

Use this with caution — it immediately releases all currently-banned IPs:

\begin{bashcode}
docker compose exec redis redis-cli --scan --pattern "ban:*" | \
  xargs docker compose exec redis redis-cli DEL
\end{bashcode}

\section{Managing Rate-Limit Windows}
\index{rate limiting!clearing}

Rate-limit sliding windows are stored as Redis Sorted Sets. To clear rate-limit
state for a specific IP (for example, after a customer reports being rate-limited):

\begin{bashcode}
# Clear by-IP rate limit for a specific address
docker compose exec redis redis-cli DEL "rl:ip:192.0.2.1"

# Clear by-IP+JA4 pair rate limit
docker compose exec redis redis-cli DEL "rl:pair:192.0.2.1:t13d1113h2_8daaf6152771_b0da82dd1658"
\end{bashcode}

To clear all rate-limit windows (useful between load tests):

\begin{bashcode}
docker compose exec redis redis-cli --scan --pattern "rl:*" | \
  xargs docker compose exec redis redis-cli DEL
\end{bashcode}

\section{Updating the GeoIP Database}
\index{GeoIP!updating}
\index{IP2Location}

The IP2Location database should be refreshed monthly. IP allocations change;
a stale database will misclassify newer IP ranges.

\begin{bashcode}
make update-geoip
\end{bashcode}

The proxy reloads the database on the next \code{SIGHUP} or can be made to
reload it immediately:

\begin{bashcode}
# Reload everything including the GeoIP database
kill -HUP $(docker compose exec ja4proxy cat /var/run/proxy.pid)
\end{bashcode}

If the database file is missing or corrupted, the proxy falls back to treating
all IPs as country code \code{XX} (unknown). Country filtering is effectively
disabled in this state, and the proxy logs a startup warning. Other functions
are unaffected.

\section{Flushing Redis Between Test Runs}
\index{Redis!flush}

When running load tests or switching between test scenarios, flush Redis to
start with a clean state:

\begin{bashcode}
docker compose exec redis redis-cli FLUSHALL
\end{bashcode}

\begin{warningbox}[FLUSHALL in production]
Never run \code{FLUSHALL} on a production Redis instance. It removes all data
including bans and enrichment caches. Use it only in development or test environments.
\end{warningbox}

\section{Checking Proxy Health}
\index{health check}

\subsection{Quick Health Check}

\begin{bashcode}
curl -s http://localhost:8080/health
\end{bashcode}

Returns JSON with status, Redis connectivity, GeoIP status, and uptime.

\subsection{Readiness Check}

\begin{bashcode}
curl -s http://localhost:8080/ready
\end{bashcode}

Returns 200 when the proxy is ready to accept connections. Used by Kubernetes
readiness probes.

\subsection{Container Status}

\begin{bashcode}
docker compose ps
docker compose top
\end{bashcode}

\subsection{Resource Usage}

\begin{bashcode}
docker stats --no-stream
\end{bashcode}

Shows CPU, memory, and network I/O for all containers. The proxy container should
be using well under 100\% CPU at normal load; spikes above this indicate a
performance issue.

\section{Log Management}
\index{logs!management}
\index{Loki}

Logs are collected by Promtail and stored in Loki. The Grafana Explore interface
queries them directly. For direct log access:

\begin{bashcode}
# Live proxy logs
docker compose logs -f ja4proxy

# Last 100 lines of all services
docker compose logs --tail=100

# Logs since a specific time
docker compose logs --since=2026-04-04T10:00:00 ja4proxy
\end{bashcode}

Log files are also written to \code{logs/} on the host. Default retention is
7 days; adjust in \code{config/proxy.yml} under \configkey{logging.retention\_days}.

\section{Scheduled Maintenance Checklist}

\begin{table}[ht]
  \small
  \begin{tabularx}{\linewidth}{llX}
    \toprule
    \textbf{Frequency} & \textbf{Task} & \textbf{Command} \\
    \midrule
    Daily & Review Grafana dashboard & Visual check: action distribution, score histogram \\
    Daily & Check for new fingerprints & Grafana Top Fingerprints panel \\
    Weekly & Review ban log & Check for patterns suggesting targeting \\
    Monthly & Update GeoIP database & \code{make update-geoip} \\
    Monthly & Review blacklist & Remove stale entries, add newly identified tools \\
    Monthly & Check Docker image updates & \code{git pull \&\& make rebuild} \\
    \bottomrule
  \end{tabularx}
  \caption{Recommended maintenance schedule}
\end{table}
